Energy consumption modeling is fundamental for ensuring the safety, reliability, and efficiency of Unmanned Aerial Vehicle (UAV) operations \cite{ren2025, muli2022, zhang2021, dorling2017}. Because drones have limited battery capacity and short endurance, often limited to 20-30 minutes, accurately predicting power demand is vital for mission planning and range estimation \cite{ren2025, muli2022, conte2022, rodrigues2021}. The sources identify four primary categories of factors that influence a drone's energy usage: drone design (e.g., weight, battery efficiency, rotor size), operating environment (e.g., wind conditions, air density, temperature), drone dynamics (e.g., flight speed, acceleration, angle of attack), and delivery operations (e.g., payload weight) \cite{ren2025, muli2022, zhang2021, patterson2018}.

\subsection{State of the art}
The methodological landscape of energy modeling is generally divided into two paradigms. First, theoretical/physics-based models, these utilize fundamental principles of flight to calculate the energy required to overcome forces such as gravity and aerodynamic drag \cite{zhang2021, rodrigues2021}. These range from simple integrated models based on lift-to-drag ratios to complex component models that separately account for induced power (to maintain lift), parasite power (to overcome air resistance) and profile power (to overcome rotational drag on propeller blades) \cite{bruehl2021,zhang2021}.

Second, data-driven/Machine Learning Models, these leverage architectures such as Long Short-Term Memory (LSTM), Bi-LSTM and Decision Trees to learn complex temporal patterns from empirical flight data \cite{ren2025, muli2022, gora2022, fatima2022}. According to the sources, these "black-box" approaches can often outperform traditional mathematical models by capturing non-linearities and environmental uncertainties without requiring detailed knowledge of a drone's internal construction \cite{muli2022, gora2022}.

A significant challenge highlighted in the sources is the lack of consensus on a universal model \cite{muli2022, zhang2021}. Comparative studies show that different modeling approaches can produce energy consumption estimates that vary by a factor of 3 to 5, even when using identical input parameters and trajectories \cite{ren2025, muli2022, rodrigues2021}. Furthermore, accurate modeling requires analyzing distinct flight regimes, as energy intensive maneuvers like vertical takeoff can account for a substantial portion of total trip energy, sometimes as much as 36\%—compared to horizontal cruise flight \cite{rodrigues2022}. Therefore, high-fidelity modeling must integrate real-world environmental data, such as wind vectors and actual battery voltage drops, to provide the predictive precision necessary for modern UAVs.

\subsection{Structure of the document}
The work begins in section \ref{data_analysis} with a detailed account of data acquisition and analysis, leveraging a publicly available dataset from 209 autonomous flights of a DJI Matrice 100 quadcopter, collected over seven months under diverse operational and environmental conditions \cite{rodrigues2021}. Building on this, the section \ref{flight_phases} defines flight phases through a threshold-based classification scheme, using statistical analysis of vertical velocity, horizontal speed, and power to distinguish between Taxi, Hover, Cruise, Climb, Descent, and Other states, ensuring physically meaningful segmentation of flight data. A physics-based power model is then developed in section \ref{Physic_model} using actuator disk theory and phase-dependent efficiency estimates, providing a transparent, first-principles framework for predicting electrical power consumption during steady-state flight. Complementing this, a Random Forest regression model is introduced in section \ref{data_driven} as a data-driven alternative, trained on 19 flight-state features—including velocity, wind, attitude, and flight mode indicators—to capture complex, non-linear relationships in energy demand. The two models are rigorously validated and compared in section \ref{task_04} using standard error metrics across all flight phases, revealing the superior accuracy of the data-driven approach while also highlighting its limitations in extrapolation and interpretability as described in section \ref{task_05}. 
